\documentclass[a4paper,10pt]{article}

\usepackage[a4paper,margin=1.15cm]{geometry}
\usepackage[T1]{fontenc}
\usepackage[utf8]{inputenc}
\usepackage[ngerman,english]{babel}
\usepackage{lmodern}
\usepackage{microtype}
\usepackage[table]{xcolor}
\usepackage{array}
\usepackage{tabularx}
\usepackage{titlesec}
\usepackage[hidelinks]{hyperref}

\definecolor{HeaderBlueDark}{RGB}{70,110,170}
\definecolor{RowBlueLight}{RGB}{235,243,255}
\definecolor{RowBlueDark}{RGB}{220,233,250}
\definecolor{SoftGray}{RGB}{90,90,90}

\setlength{\parindent}{0pt}
\setlength{\tabcolsep}{5pt}
\renewcommand{\arraystretch}{1.22}
\newcolumntype{Y}{>{\centering\arraybackslash}X}
\newcommand{\TableHeader}[1]{\color{white}\textbf{#1}}

\titleformat{\section}[block]
{\Large\bfseries\color{HeaderBlueDark}}
{}{0pt}{}
\titlespacing{\section}{0pt}{0.2em}{0.35em}

\titleformat{\subsection}[block]
{\normalsize\bfseries\color{HeaderBlueDark}}
{}{0pt}{}
\titlespacing{\subsection}{0pt}{0.75em}{0.25em}

\title{\vspace{-1.4cm}\textcolor{HeaderBlueDark}{Artikel in den vier Kasus}\\[-1mm]
{\large\color{SoftGray}Articles in the four cases}}
\date{}

\begin{document}
\maketitle
\vspace{-1.25cm}

\section{Datentabelle / Data table}

\subsection{Bestimmter Artikel / Definite article}
\rowcolors{2}{RowBlueLight}{RowBlueDark}
\begin{tabularx}{\textwidth}{>{\bfseries}p{3.1cm} YYYY}
\rowcolor{HeaderBlueDark}
\TableHeader{Kasus / Case} & \TableHeader{Maskulin} & \TableHeader{Feminin} & \TableHeader{Neutrum} & \TableHeader{Plural} \\
Nominativ & der & die & das & die \\
Akkusativ & den & die & das & die \\
Dativ & dem & der & dem & den \\
Genitiv & des & der & des & der \\
\end{tabularx}

\subsection{Unbestimmter Artikel / Indefinite article}
\rowcolors{2}{RowBlueLight}{RowBlueDark}
\begin{tabularx}{\textwidth}{>{\bfseries}p{3.1cm} YYYY}
\rowcolor{HeaderBlueDark}
\TableHeader{Kasus / Case} & \TableHeader{Maskulin} & \TableHeader{Feminin} & \TableHeader{Neutrum} & \TableHeader{Plural} \\
Nominativ & ein & eine & ein & -- \\
Akkusativ & einen & eine & ein & -- \\
Dativ & einem & einer & einem & -- \\
Genitiv & eines & einer & eines & -- \\
\end{tabularx}

\subsection{Negativartikel / Negative article}
\rowcolors{2}{RowBlueLight}{RowBlueDark}
\begin{tabularx}{\textwidth}{>{\bfseries}p{3.1cm} YYYY}
\rowcolor{HeaderBlueDark}
\TableHeader{Kasus / Case} & \TableHeader{Maskulin} & \TableHeader{Feminin} & \TableHeader{Neutrum} & \TableHeader{Plural} \\
Nominativ & kein & keine & kein & keine \\
Akkusativ & keinen & keine & kein & keine \\
Dativ & keinem & keiner & keinem & keinen \\
Genitiv & keines & keiner & keines & keiner \\
\end{tabularx}



\subsection{jeder/jede/jedes-Familie / every-each family}
\rowcolors{2}{RowBlueLight}{RowBlueDark}
\begin{tabularx}{\textwidth}{>{\bfseries}p{3.1cm} YYYY}
\rowcolor{HeaderBlueDark}
\TableHeader{Kasus / Case} & \TableHeader{Maskulin} & \TableHeader{Feminin} & \TableHeader{Neutrum} & \TableHeader{Plural} \\
Nominativ & jeder & jede & jedes & -- \\
Akkusativ & jeden & jede & jedes & -- \\
Dativ & jedem & jeder & jedem & -- \\
Genitiv & jedes & jeder & jedes & -- \\
\end{tabularx}

\vspace{0.35em}
{\small\color{SoftGray}\textbf{Grammatische Rolle / Grammatical role:} Vor einem Nomen ist \textbf{jeder/jede/jedes} ein indefinites Artikelwort beziehungsweise Determinativ. Allein benutzt ist es ein Indefinitpronomen. / Before a noun, \textbf{jeder/jede/jedes} is an indefinite article word or determiner. Used alone, it is an indefinite pronoun.}

\vspace{0.35em}
{\small\color{SoftGray}\textbf{Hinweis / Note:} \textbf{jeder/jede/jedes} wird normalerweise mit Singular benutzt. Fuer Plural benutzt man meistens \textbf{alle}. / \textbf{jeder/jede/jedes} is normally used with singular. For plural, German usually uses \textbf{alle}.}

\vspace{0.5em}
{\small\color{SoftGray}\textbf{Hinweis / Note:} Der unbestimmte Artikel hat keinen Plural. / The indefinite article has no plural form.}

\end{document}
